% EM structure and teaching revision 20260928
\subsection{机电耦合：两条方程与一份能量账}
导体棒长 $\ell$、质量 $m$，在匀强磁场中沿导轨运动。忽略回路自感、取总电阻 $R$，则感应电流大小为 $I=B\ell v/R$，磁力反向阻碍运动，故
\[
 m\dot v=F_{\mathrm{ext}}-\frac{B^2\ell^2}{R}v.
\]
恒定外力下，终端速度 $v_\infty=F_{\mathrm{ext}}R/(B^2\ell^2)$，趋近它的时间常数为 $\tau_m=mR/(B^2\ell^2)$。当回路电感 $L$ 不可忽略且近似不随棒的位置变化时（例如固定外接电感占主导），取电感电压关系 $u_L=L\dot I$，必须同时求解
\[
 L\dot I+RI=B\ell v,\qquad m\dot v=F_{\mathrm{ext}}-B\ell I.
\]
若 $L/R\ll\tau_m$，电流才可近似及时跟随速度，退回第一式。把两式分别乘以 $I$ 和 $v$ 后相加，得到
\[
 F_{\mathrm{ext}}v=\frac{\dd}{\dd t}\left(\frac12mv^2+\frac12LI^2\right)+I^2R.
\]
自感电压关系及其能量将在本章最后一节推导。若运动使 $L(x)$ 显著变化，还会产生相应的电动势和自感磁力，不能继续使用这里的常量 $L$ 模型。这条能量式能检验耦合项符号：机械能转出的 $B\ell Iv$ 必须恰好进入电路，不能在两边同时作为损耗。接入电容、电源或增加第二根可动棒时，仍按同样顺序增列各自的状态变量与约束。

